This paper is about one narrow, common regime: post-training a generative
policy that was pretrained on a small set of demonstrations, online, under
sparse binary reward. Our experiments locate the boundary of that regime
from both sides. Inside it, any update that lets the critic displace actions
without an action-space bound destroys the policy --- in our system and in
the official implementations of an entire method family --- because the
critic is only competent in the narrow tube its data covers and sparse
reward offers no route back once the policy leaves it. Outside it, the claim weakens but does not vanish: under dense reward the
failure stops being terminal, yet on three DMC tasks \method{} still wins
--- dense reward removes the catastrophe, not the cost of wandering, and
the unconstrained update still falls below its base on quadruped
(\cref{tab:dmc-scope}). Only at saturating critic coverage
(\cref{fig:toy}a) does backpropagation become the better tool. Within the regime, the
missing update rule is field-target regression under an action-space bound
anchored to the frozen base. \method{} expresses reward and anchoring as composable velocity fields over
action particles, and recovers the backpropagated direction in the
small-step limit. Its working component is the action-space trust region:
the critic sets the direction, while a bound --- a pointwise clip on each
particle's displacement, a restoring pull outside a radius of the frozen
base, or both --- sets how far any single update may carry the policy.

Given the \emph{same} base, critic, data and budget, the field-target form
extracts far more than backpropagating that critic: $+55$ points on
robomimic square, and $+9.4$ versus $+0.9$ across the eight hardest
LIBERO-90 tasks, with the same ordering after swapping in a stronger
flow-matching base --- so the effect belongs to the update rule, not the
starting point. Against a critic-free policy gradient at matched budget,
\method{} is at parity while needing none of the noise injection that makes
likelihood ratios definable for a one-step generator. The ablations then
pin the mechanism: removing every action-space bound gives exactly $0.000$
success on all five factorial tasks, the step clip alone also gives
$0.000$, the always-on penalty ($\rho{=}0$, exactly a BC term) falls back
to base, and only configurations with a restoring term toward the frozen
base survive. The constraint that matters is not a limit on how fast the
policy moves but a tether to where it started.

\textbf{Limitations.}
\label{sec:limitations}
Three specific results run against our thesis and we state them
plainly. On the Study-3 caddy task the backpropagated actor ends above \method{}
($0.653$ versus $0.635$); on two flow-matching tasks the
\emph{un-finetuned} base beats every fine-tuner including ours ($0.917$
on the orange-juice task, $0.620$ on the right-caddy task); and on robomimic can the flow-matching pipeline
edges us on two of three seeds. Single-seed variation is also
substantial on LIBERO --- repeating the frozen recipe on the red-mug task with a
second seed moves the result by $8.9$ points --- so per-task deltas
smaller than roughly ten points, including our $\eta_Q$ and
$\delta_{\mathrm{tot}}$ sweeps, indicate direction rather than magnitude.
The claims that survive this variation are the ones separated by margins
far larger than it: the $+9.4$ versus $+0.9$ mean gap, the
non-overlapping step-size bands, and the below-base collapse when the
critic's gradient is deleted.

The gains track how much a critic actually knows.
When a single critic is shared across all eight LIBERO tasks its
action-gradient collapses and \emph{no} critic-based update --- ours or the
state of the art's --- improves on its base; when a base is already near its
ceiling, \method{} can end slightly below it ($-2.0$ and $-1.2$ on two
flow-matching tasks), since our trust region is calibrated for lifting weak
bases rather than for refining strong ones. Our critic-trust diagnosis is
also post hoc: choosing the reward-field resolution online from critic-trust
signals is the natural next step, as is testing whether the reward step
should adapt per task. Quantitative fine-tuning results are currently limited
to simulation; \cref{sec:exp-real-robot} describes the ongoing physical-robot
study, whose four-way success-rate comparison remains to be completed.

We should be precise about which part of the story generalizes, because two
claims are easy to conflate. The \emph{update rule} is base-agnostic: it
consumes a frozen noise-to-action map, a residual head and a critic, and
never touches how the base draws its action, and it transfers unchanged to a
flow-matching base, where dropping it into DICE-RL raises success on square
from $0.831$ to $0.934$; the diffusion-base pair converges to a shared
${\approx}0.96$ ceiling with \method{} arriving $1.6\times$ sooner
(\cref{tab:baselines}). Its \emph{size} of benefit, however, is not base-agnostic: on the
flow-matching LIBERO base the same recipe is a wash against the base
($0.756$ vs $0.753$) and only $+1.3$ points over DICE-RL across the seven
tasks where both ran. What does \emph{not} generalize is the
economy of the construction. A drifting model is pretrained by
field-target regression, so on that base the fine-tuner reuses machinery
the policy class already has; a flow-matching model is pretrained by
conditional velocity regression, so there the same update is bolted on
rather than inherited (and costs one base forward pass per particle per
integration step). The principle we can defend is therefore narrow: \emph{what keeps a
critic-driven update on the data is a bound on the total displacement in
action space, not a limit on step size}. The claim stops at the \emph{location} of the constraint. On the two LIBERO tasks a hinge
penalty and a hard projection onto $\|r\|\le\rho$ both beat our dead-zone
anchor, while on square the anchor wins by a wide margin
(\cref{tab:locus}); the anchor is one workable instrument rather than the
necessary one, and which instrument is best appears task-dependent. The distinction is not cosmetic. A step limit constrains
speed, and a drift assembled from permitted small steps still arrives; in our
isolation the limit alone leaves the policy at zero success. A
behaviour-cloning penalty restores, but does so everywhere, including the
region where the critic is reliable, and must therefore trade far-field
containment against near-field progress with a single coefficient; it also
fails alone. The dead zone separates those jobs.

We deliberately claim no more than this. We have no principled rule for
choosing $\rho$ --- it is tuned, not derived --- and the argument for why the
dead zone should exist is an argument about matching the regularizer's
spatial profile to the critic's reliability, not a bound. Deriving $\rho$
from a measure of local critic reliability is the obvious next question, and
we leave it open.
